 
\def\filename{leaflet-manual.tex}
\def\fileversion{v1.0d}   % change this when leaflet-manual changed, too.
\def\filedate{2012/06/04}
\def\docdate {2012/06/04} % change this when leaflet-manual changed, too.
\listfiles
\errorcontextlines=99
\documentclass[
notumble,
nofoldmark,
%%dvipdfm,
%%portrait,
%%titlepage,
%%nocombine,
%%a3paper,
%%debug,
%%nospecialtricks,
%%draft,
]{leaflet}
\usepackage{tabularx}
\usepackage{lipsum}
\usepackage{eurosym}
\renewcommand*\foldmarkrule{.3mm}
\renewcommand*\foldmarklength{5mm}

\usepackage[T1]{fontenc}
% \usepackage[T3]{inputenc}
\usepackage{textcomp}
\usepackage{mathptmx}
\usepackage{libertine}
\usepackage[scaled=0.9]{helvet}
\makeatletter
\def\ptmTeX{T\kern-.1667em\lower.5ex\hbox{E}\kern-.075emX\@}
\DeclareRobustCommand{\ptmLaTeX}{L\kern-.3em
        {\setbox0\hbox{T}%
         %\vb@xt@ % :-)
         \vbox to\ht0{\hbox{%
                            \csname S@\f@size\endcsname
                            \fontsize\sf@size\z@
                            \math@fontsfalse\selectfont
                            A}%
                      \vss}%
        }%
        \kern-.12em
        \ptmTeX}
\makeatother
\let\TeX=\ptmTeX
\let\LaTeX=\ptmLaTeX
\usepackage{shortvrb}
\MakeShortVerb{\|}
\usepackage{url}
\usepackage{graphicx}
\usepackage[dvipsnames,usenames]{color}
\definecolor{LIGHTGRAY}{gray}{.9}
\definecolor{lspblue1}{RGB}{208, 236, 246} 
\definecolor{lsptmnlp}{cmyk}{0.9,0.5,0.15,0.3}
% \usepackage{langscicolors} 
\newcommand\Lpack[1]{\textsf{#1}}
\newcommand\Lclass[1]{\textsf{#1}}
\newcommand\Lopt[1]{\texttt{#1}}
\newcommand\Lprog[1]{\textit{#1}}

\newcommand*\defaultmarker{\textsuperscript\textasteriskcentered}

 

\CutLine*{1}% Dotted line without scissors
\CutLine*{3}%   
\CutLine*{4}%   
\CutLine*{6}%   

\AddToBackground*{1}{%  Background of a small page
  \put(0,0){\textcolor{lsptmnlp}{\rule{\paperwidth}{\paperheight}}}}

\AddToBackground*{2}{%  Background of a small page
  \put(0,0){\textcolor{lsptmnlp}{\rule{\paperwidth}{\paperheight}}}}
 
\AddToBackground{2}{
    \put(-2,0){
        \includegraphics[width=.3\paperwidth]{lsp-rainbow-large.png}%
    } 
}

\AddToBackground{3}{
    \put(-2,0){
        \includegraphics[width=.3\paperwidth]{lsp-rainbow-large.png}%
    } 
}

\AddToBackground{4}{
    \put(-2,0){
        \includegraphics[width=.3\paperwidth]{lsp-rainbow-large.png}%
    } 
}


\begin{document}
\newsavebox{\pageone}
\newsavebox{\pagetwo}

\savebox{\pageone}{
    \parbox{\textwidth}{
    \color{LIGHTGRAY} 
    \thispagestyle{empty}
    
    \noindent 
    ~
    \parbox{.088\paperwidth}{
      \color{LIGHTGRAY}
      \vspace*{-.5cm}
      \includegraphics[width=.3\paperwidth]{langsci_logo_nocolor.pdf} 
    }
    \parbox{.8\paperwidth}{
      \raggedleft 
      \vspace*{-2.5cm}{\includegraphics[width=.4\paperwidth]{lsp-rainbow-large.png}}
    }

    \vspace*{11cm}

    \raggedleft{
    \mbox{\Huge Language Science Press}
    {\LARGE Translation and Multilingual Natural Language Processing}}
    \raggedright 

    \medskip 
\small 
The book series centres around human and machine translation, with a special emphasis on empirical studies. This includes computational, corpus linguistic and cognitive aspects of translation.
By its nature, the topic of translation is interdisciplinary in the sense that it involves many of the classical linguistic sub-disciplines such as computational linguistics, corpus linguistics, morphology, syntax, semantics, pragmatics, text linguistics, lexicography, psycholinguistics, neurolinguistics, applied linguistics and others. 
% However, all book submissions need to have a clear focus on the translation aspect, and a special emphasis is laid on empirical studies. The aim of the book series is to bring these different perspectives closer together by offering a forum for all different approaches to the empirical study of translation. The series welcomes in particular studies investigating corpus data and/or experimental findings, preferably in – but not limited to -- a quantitative perspective. 


  }
}

\savebox{\pagetwo}{
  \parbox{\textwidth}{
    \color{LIGHTGRAY}
    \section{\color{LIGHTGRAY} \sffamily \Large Possible Topics} 
 
\begin{itemize}
 \item [$\rangle$]machine translation (statistical, rule-based, example-based and hybrid), machine aided translation and interpreting
translation technology\item [$\rangle$]
use of corpora in translation\item [$\rangle$]
annotation, alignment and searchability of translation data\item [$\rangle$]
cognitive aspects of translation\item [$\rangle$]
modelling the translation process\item [$\rangle$]
multi-modal and audiovisual translation 
\end{itemize}



    \section{\color{LIGHTGRAY} \sffamily \Large Editorial Board}    

    \begin{itemize}
    \item[$\rangle$] Reinhard Rapp (Chief Editor, Aix-Marseille Universit\'e)
    \item[$\rangle$] Silvia Hansen-Schirra (Johannes Gutenberg-Universit\"at Mainz)
    \item[$\rangle$] Oliver \v{C}ulo (Johannes Gutenberg-Universit\"at Mainz)
    \end{itemize}
    
    \section{\color{LIGHTGRAY} \sffamily \Large Support us!}
	\begin{itemize}
	  \item[$\rangle$] submit your manuscript 
	  \item[$\rangle$] become a community proofreader 
	  \item[$\rangle$] become a community typesetter
	  \item[$\rangle$] become a public supporter
	  \item[$\rangle$] donate
	\end{itemize} 
    \section{\color{LIGHTGRAY} \sffamily \Large Contact} 
	\begin{itemize}
	  \item[$\rangle$] \sffamily www.langsci-press.org/catalog/series/tmnlp
	  \item[$\rangle$]\sffamily  tmnlp@langsci-press.org 
	\end{itemize} 
  }
}

\usebox{\pageone}

\newpage

\usebox{\pagetwo}

\newpage

\usebox{\pagetwo}

\newpage

\usebox{\pagetwo}

\newpage

\usebox{\pageone}

\newpage

\usebox{\pageone}
 

\loggingall
\end{document}
\endinput 
